% 中文摘要

随着在线教育的普及，慕课（MOOC）等平台为学习者提供了丰富的学习资源，但同时也面临着学习者个体差异大、难以提供个性化指导以及缺乏及时有效反馈的挑战。传统的在线学习平台多为静态的内容展示，难以适应不同学习者的背景和进度。此外，在代码实训和视频学习等核心场景中，学生往往在遇到问题时无法获得针对性的启发式辅导，导致学习效率低下。为解决上述问题，本文设计并实现了一款基于大语言模型（LLM）和检索增强生成（RAG）技术的个性化学习系统。

本文首先对系统的需求进行了全面分析，明确了基于知识库的视频问答、智能代码实训以及基于抽象语法树（AST）的代码分析等核心功能。在系统架构上，采用了前后端分离的模式，前端使用 Vue 3 框架，后端基于 Flask 搭建，并利用 MySQL 和 Redis 进行数据与状态管理。在此基础上，本文深入探讨了系统的关键技术实现。在代码实训模块，提出并实现了一种结合苏格拉底式启发策略的算法，该算法通过解析学生提交代码的 AST 结构，结合大语言模型的推理能力，为学生提供渐进式的调试指引，而非直接给出答案，从而培养学生的独立解题能力。在视频学习模块，系统引入了 RAG 技术，将课程视频的字幕和文档转换为向量存储于 FAISS 数据库中。当学习者提出问题时，系统能够在特定视频的知识库中检索相关上下文，并生成带有精确时间戳引用的回答，实现精准的答疑交互。

最后，本文对系统进行了全面的功能和性能测试。测试结果表明，该系统不仅在白盒测试中达到了 100\% 的代码覆盖率，而且在并发性能测试中表现稳定，能够满足实际教学场景下的使用需求。本系统的设计与实现，不仅验证了大模型技术在教育垂直领域的应用潜力，也为未来智能化、个性化的在线教育平台建设提供了有益的探索和实践参考。